\section{Evaluating Generated Reviews}
\label{sec:bg-evaluation}

How to score a generated review is contested. Text-overlap metrics against a
single human reference have been shown to misrank review comments, because a
change admits many valid reviews~\cite{hong2024deepcrceval}. Criterion-based
rubrics scored by a language model replaced them: Liu et al.'s G-Eval
established the pattern, reporting correlation with human ratings markedly
better than $n$-gram metrics~\cite{liu2023geval}. Zheng et al.\ generalised it
with MT-Bench, finding agreement between a strong judge and human preferences
comparable to agreement between two humans~\cite{zheng2023judging}. He et al.\
survey adoption of the paradigm in software
engineering~\cite{he2025llmasjudgese}.

The same literature catalogues failure modes. Zheng et al.\ document position
bias, verbosity bias, and self-enhancement bias~\cite{zheng2023judging};
Panickssery et al.\ show that models recognise and favour their own
generations~\cite{panickssery2024selfpreference}. Verbosity bias is
particularly relevant here: adding context to a prompt tends to make the output
longer, so a judge that rewards length would reward context augmentation for the
wrong reason.

The mitigations this thesis adopts follow directly. Using several judges from
different providers limits any one family's idiosyncrasies. Recomputing the
effect under individual judges makes evaluator sensitivity visible. Blinding
and pairwise presentation control position and identity
effects~\cite{chiang2024chatbotarena}. Where objective ground truth can be
constructed, adjudication is checked against it rather than trusted.

One conceptual limitation applies to every rubric of this kind. Asking whether
a review addresses a criterion measures \emph{coverage}: whether the review
says something about the subject. It does not measure whether what it says is
\emph{correct}. This thesis uses a coverage rubric for its main comparison and
supplements it with both an oracle-scored experiment and a human study,
precisely because coverage alone cannot answer whether added context was used
correctly.

The evaluation methods specific to code review share this gap. DeepCRCEval
replaces overlap metrics with a criterion-based rubric scored by a single
judge~\cite{hong2024deepcrceval}; CRScore constructs pseudo-references from
static analysers~\cite{naik2025crscore}. Neither reports agreement between
independent scorers, because neither has two. This matters more for a
context-augmentation study than for a model ranking, because verbosity bias
from the treatment cannot be separated from a real improvement without a panel.
